\documentclass[11pt,a4paper]{article}

\usepackage[utf8]{inputenc}
\usepackage[T1]{fontenc}
\usepackage[spanish]{babel}
\usepackage{lmodern}
\usepackage{geometry}
\usepackage{amsmath,amssymb}
\usepackage{booktabs}
\usepackage{array}
\usepackage{xcolor}
\usepackage{hyperref}
\usepackage{fancyhdr}
\usepackage{setspace}

\geometry{left=2.6cm,right=2.6cm,top=2.5cm,bottom=2.5cm}
\definecolor{azul}{HTML}{17365D}
\hypersetup{colorlinks=true,linkcolor=azul,urlcolor=azul}
\setstretch{1.13}
\setlength{\parskip}{5pt}

\pagestyle{fancy}
\fancyhf{}
\fancyhead[L]{\small Fundamentos de Inteligencia Artificial}
\fancyhead[R]{\small Yachay Tech}
\fancyfoot[C]{\thepage}
\setlength{\headheight}{14pt}

\begin{document}

% ==================================================
% PORTADA
% ==================================================

\begin{titlepage}
\centering
\vspace*{1.5cm}

{\Large\bfseries UNIVERSIDAD YACHAY TECH\par}
\vspace{0.3cm}
{\large Maestría en Ciencia de Datos\par}

\vspace{1.2cm}
\textcolor{azul}{\rule{\textwidth}{1.5pt}}

\vspace{0.8cm}
{\Huge\bfseries
Predicción de alertas de\\[0.2cm]
desnutrición infantil\par}

\vspace{0.5cm}
{\Large
Evaluación comparativa de modelos de\\
aprendizaje automático y análisis estocástico\par}

\vspace{0.8cm}
\textcolor{azul}{\rule{\textwidth}{1.5pt}}

\vspace{1.4cm}

{\large\textbf{Asignatura}\par}
{\large Fundamentos de Inteligencia Artificial\par}

\vspace{0.8cm}

{\large\textbf{Autor}\par}
{\Large\bfseries Mat. Héctor Alexander Román Potosí\par}

\vspace{0.8cm}

{\large\textbf{Docente}\par}
{\large Ph.D. Manuel Eugenio Morocho Cayamcela\par}

\vfill

{\large Examen de medio semestre -- IISEM2026\par}
{\large 31 de octubre de 2026\par}

\end{titlepage}

% ==================================================
% RESUMEN
% ==================================================

\section*{Resumen}

Este trabajo estudia un problema de clasificación binaria
de alertas de desnutrición infantil en Ecuador mediante un
conjunto de datos sintético de 70\,000 observaciones de
entrenamiento y 30\,000 observaciones de prueba.

Se evaluaron clasificadores lineales, ensambles de árboles,
Gradient Boosting, una red neuronal multicapa y enfoques
probabilísticos. Se priorizó la exactitud
(\emph{accuracy}), métrica oficial de la competencia.

Los experimentos no identificaron una mejora reproducible
frente al clasificador de clase mayoritaria, cuyo accuracy
en validación fue aproximadamente del 89.99\%.

La solución enviada a Kaggle mediante Random Forest
alcanzó un accuracy público de 91.1255\%.
Sin embargo, todas sus predicciones fueron de clase cero,
por lo que su desempeño equivale al baseline.

Se presenta una interpretación estadística de este
resultado y se discuten las limitaciones del conjunto
sintético y de la evaluación mediante accuracy.

\noindent\textbf{Palabras clave:}
clasificación binaria, Random Forest, XGBoost,
desbalance de clases, inferencia estocástica,
accuracy, Kaggle.

\tableofcontents
\newpage

% ==================================================
\section{Introducción}
% ==================================================

La identificación de alertas de desnutrición infantil
constituye una tarea de clasificación supervisada.
A partir de características individuales se busca
estimar una variable de respuesta binaria.

El presente trabajo utiliza un conjunto sintético
publicado para la competencia académica
\emph{Ecuador Infant Malnutrition Competition
(IISEM2026 MCD)}.

Es importante distinguir la naturaleza académica de
estos datos de una herramienta clínica:
las etiquetas sintéticas no equivalen a diagnósticos
médicos verificados.

\subsection{Objetivo general}

Comparar diversos modelos de aprendizaje automático
para maximizar la exactitud de clasificación sobre
datos no observados, aplicando procedimientos
reproducibles de entrenamiento y validación.

\subsection{Objetivos específicos}

\begin{itemize}
\item Examinar la estructura y calidad de los datos.
\item Construir un baseline de clase mayoritaria.
\item Comparar algoritmos de distintas familias.
\item Evaluar el efecto del desbalance y los umbrales.
\item Estudiar una explicación probabilística
      del rendimiento observado.
\item Generar una entrega válida para Kaggle.
\end{itemize}

% ==================================================
\section{Descripción de los datos}
% ==================================================

La base de entrenamiento contiene 70\,000 registros
y seis variables predictoras:

\begin{center}
\begin{tabular}{ll}
\toprule
Variable & Descripción\\
\midrule
\texttt{edad\_meses} & Edad en meses\\
\texttt{sexo} & Categoría de sexo\\
\texttt{region} & Región geográfica\\
\texttt{nivel\_educacion\_madre} & Escolaridad materna\\
\texttt{peso\_kg} & Peso en kilogramos\\
\texttt{talla\_cm} & Talla en centímetros\\
\bottomrule
\end{tabular}
\end{center}

La variable objetivo es
\texttt{alerta\_desnutricion}, con valores cero y uno.

\subsection{Distribución de clases}

Se observaron 62\,991 registros de clase cero
y 7\,009 de clase uno:

\[
P(Y=0)\approx0.899871,
\qquad
P(Y=1)\approx0.100129.
\]

Existe, por tanto, un importante desbalance de clases.

\subsection{Calidad de los datos}

Se identificaron valores faltantes en
\texttt{peso\_kg} (aproximadamente 4.96\%) y
\texttt{talla\_cm} (aproximadamente 10.07\%).

También se observaron 3\,921 edades negativas
y 3\,163 edades superiores a 24 meses.

Estas anomalías se analizaron sin eliminar
arbitrariamente observaciones etiquetadas.

% ==================================================
\section{Metodología}
% ==================================================

\subsection{Preparación}

Se utilizó imputación de mediana en las variables
numéricas y codificación One-Hot para las categóricas.

El preprocesamiento se incorporó dentro de
\texttt{Pipeline} y \texttt{ColumnTransformer}
para impedir que los estadísticos de validación
intervinieran en el ajuste de entrenamiento.

\subsection{Métrica principal}

La métrica oficial es:

\[
\operatorname{Accuracy}
=
\frac{TP+TN}{TP+TN+FP+FN}.
\]

Se calcularon además recall, F1, ROC-AUC y
balanced accuracy como métricas diagnósticas.

\subsection{Modelo de referencia}

El clasificador de clase mayoritaria predice:

\[
\widehat{Y}(x)=0.
\]

Su exactitud esperada se aproxima a la proporción
de observaciones pertenecientes a la clase cero.

\subsection{Modelos estudiados}

Se experimentó con DummyClassifier, regresión
logística, Random Forest, Extra Trees,
HistGradientBoosting, LightGBM, CatBoost,
XGBoost y una red neuronal MLP.

Adicionalmente se analizaron probabilidades
beta-binomiales y un modelo estocástico Probit.

% ==================================================
\section{Resultados experimentales}
% ==================================================

\subsection{Validación cruzada}

En tres particiones estratificadas se obtuvieron:

\begin{center}
\small
\begin{tabular}{lrrr}
\toprule
Modelo & Accuracy & Recall & ROC-AUC\\
\midrule
Dummy & 0.89987 & 0.00000 & 0.50000\\
Random Forest & 0.89987 & 0.00000 & 0.65429\\
HistGradientBoosting & 0.89987 & 0.00000 & 0.65474\\
RF balanceado & 0.45807 & 0.85633 & 0.65305\\
HGB balanceado & 0.42281 & 0.92338 & 0.65003\\
\bottomrule
\end{tabular}
\end{center}

Los modelos balanceados incrementaron el recall
de la clase positiva, pero redujeron notablemente
la exactitud, por lo que no fueron seleccionados
para la métrica de la competencia.

\subsection{Comparación competitiva}

En un torneo adicional se evaluaron Random Forest,
Extra Trees, LightGBM y CatBoost.

Ninguna configuración superó el accuracy del
clasificador mayoritario.

La auditoría de las predicciones con mayor riesgo
tampoco encontró grupos con una precisión
suficientemente alta para justificar el cambio
de cero a uno.

\subsection{XGBoost}

Se realizó una evaluación mediante cinco
particiones estratificadas.

El accuracy global fuera de muestra por fold
(\emph{out-of-fold}) fue:

\[
A_{\mathrm{XGB}}=0.899871.
\]

El ROC-AUC global fue aproximadamente 0.653106.

Los cinco folds predijeron exclusivamente
la clase cero con el umbral 0.5.

\subsection{Red neuronal}

Una red neuronal MLP con capas ocultas
de 32 y 16 neuronas obtuvo ROC-AUC aproximado
de 0.599958, sin mejorar el accuracy
de la clase mayoritaria.

% ==================================================
\section{Análisis matemático y estocástico}
% ==================================================

\subsection{Regla bayesiana}

Para una clasificación binaria con igual costo
para ambos tipos de error, la regla óptima
según accuracy es:

\[
\widehat y(x)=
\begin{cases}
1,&P(Y=1\mid X=x)>0.5,\\
0,&P(Y=1\mid X=x)\leq0.5.
\end{cases}
\]

Por tanto, un modelo puede discriminar probabilidades
sin que ello produzca predicciones positivas.

\subsection{Modelo beta-binomial}

Se estimaron proporciones condicionales de alerta
por región y escolaridad materna mediante
una distribución beta-binomial.

El modelo obtuvo ROC-AUC aproximado de 0.630545
y Brier Score de 0.087296.

Ninguno de los grupos analizados presentó
una probabilidad posterior media cercana a 0.5.

\subsection{Modelo Probit}

Para estudiar el carácter estocástico de las
etiquetas, se planteó:

\[
Z=r(X)+\varepsilon,
\qquad
\varepsilon\sim\mathcal N(0,1),
\]

con regla de etiquetado determinada por un umbral
\(c\). En esa aproximación:

\[
P(Y=1\mid X=x)=\Phi(r(x)-c),
\]

donde \(\Phi\) es la función de distribución
acumulada normal estándar.

El experimento estimó:

\[
\widehat c=1.065837,
\qquad
\max_x\widehat P(Y=1\mid X=x)=0.143249.
\]

El Probit obtuvo ROC-AUC 0.649634 y
Brier Score 0.086889.

Como la función de riesgo reconstruida cumple
\(r(x)\leq0\), este modelo estima probabilidades
inferiores a 0.5 para todas las observaciones.

\textbf{Limitación:} la reconstrucción del mecanismo
generador es aproximada. Este resultado no prueba
que ningún algoritmo imaginable pueda superar
el baseline sobre cualquier muestra.

% ==================================================
\section{Comprobación de generalización}
% ==================================================

Se compararon las distribuciones de las variables
predictoras entre entrenamiento y prueba.

Un clasificador entrenado para distinguir
ambos conjuntos obtuvo:

\[
\mathrm{ROC\text{-}AUC}_{\mathrm{origen}}
=0.495746.
\]

No se detectó una separación importante
entre las características observadas de los dos
conjuntos.

La estimación beta-binomial del accuracy de
predecir siempre cero sobre el test completo fue:

\[
\widehat A_{\mathrm{test}}=0.899172.
\]

El intervalo predictivo aproximado al 95\%
fue de 89.5161\% a 90.3182\%.

Este intervalo depende de los supuestos del
modelo utilizado y no constituye garantía
de la puntuación final.

% ==================================================
\section{Resultado de Kaggle}
% ==================================================

El modelo Random Forest se entrenó sobre las
70\,000 observaciones etiquetadas y generó
30\,000 predicciones.

El archivo entregado cumplió la estructura:

\begin{center}
\texttt{ID, alerta\_desnutricion}
\end{center}

Las 30\,000 predicciones fueron cero.

\subsection{Puntuación pública}

La plataforma registró el envío con
referencia \texttt{57815647} y obtuvo:

\[
\boxed{
\mathrm{Accuracy}_{\mathrm{publico}}
=0.911254675280
}
\]

Esto equivale a aproximadamente 91.1255\%
de exactitud sobre la fracción pública.

\textbf{Interpretación:} dado que las predicciones
fueron todas cero, esta puntuación es exactamente
la del baseline constante y no evidencia una
mejora de Random Forest frente a DummyClassifier.

La puntuación privada no estaba disponible
en el momento de documentar el resultado.

% ==================================================
\section{Discusión}
% ==================================================

Los resultados muestran que aumentar la
complejidad del modelo no produjo ganancias
de accuracy bajo los protocolos evaluados.

El fuerte desbalance de clases y el mecanismo
estocástico sintético proporcionan una
explicación coherente para esta observación.

El ROC-AUC superior a 0.5 en diversos modelos
indica cierta capacidad de ordenación de riesgos,
pero no es suficiente para obtener más aciertos
al aplicar el umbral convencional.

El resultado de 91.1255\% corresponde a
una métrica de competencia sobre datos sintéticos.
No debe interpretarse como una validación
de utilidad clínica del modelo.

% ==================================================
\section{Conclusiones}
% ==================================================

\begin{enumerate}
\item Se construyó un flujo reproducible
de clasificación supervisada.
\item Se analizaron problemas de calidad,
desbalance y posibles anomalías.
\item Se compararon distintas familias de modelos,
incluyendo aprendizaje profundo y enfoques
probabilísticos.
\item Ningún modelo demostró una mejora consistente
de accuracy frente a la clase mayoritaria.
\item El análisis estocástico permitió formular
una explicación matemática del comportamiento
de los clasificadores.
\item Se generó y validó una submission oficial
con accuracy público de 91.1255\%.
\end{enumerate}

\section*{Reproducibilidad}

Se adjuntan el notebook ejecutable, los scripts
de experimentación, los resultados principales,
los archivos oficiales de entrada y la evidencia
de la participación en Kaggle.

No se emplearon etiquetas ocultas del test
para ajustar, seleccionar o evaluar los modelos.

% ==================================================
\begin{thebibliography}{9}

\bibitem{breiman}
Breiman, L. (2001).
\emph{Random Forests}.
Machine Learning, 45, 5--32.

\bibitem{grinsztajn}
Grinsztajn, L., Oyallon, E. y Varoquaux, G.
(2022).
\emph{Why do tree-based models still outperform
deep learning on typical tabular data?}
NeurIPS.

\bibitem{mcel}
McElfresh, D. y colaboradores (2023).
\emph{When Do Neural Nets Outperform
Boosted Trees on Tabular Data?}
NeurIPS.

\bibitem{repo}
Morocho Cayamcela, M. E.
\emph{ArtificialIntelligence: repositorio
académico y generador de datos sintéticos}.
\url{https://github.com/eugeniomorocho/ArtificialIntelligence}

\end{thebibliography}

\end{document}
